% Intended LaTeX compiler: pdflatex
\documentclass[a4paper,12pt]{article}
\usepackage[utf8]{inputenc}
\usepackage[T1]{fontenc}
\usepackage{graphicx}
\usepackage{longtable}
\usepackage{wrapfig}
\usepackage{rotating}
\usepackage[normalem]{ulem}
\usepackage{amsmath}
\usepackage{amssymb}
\usepackage{capt-of}
\usepackage{hyperref}
\usepackage{\string~/.emacs.d/common}
\author{mahmood}
\date{\today}
\title{decision tree}
\hypersetup{
 pdfauthor={mahmood},
 pdftitle={decision tree},
 pdfkeywords={},
 pdfsubject={},
 pdfcreator={Emacs 28.2 (Org mode 9.5.5)}, 
 pdflang={English}}
\begin{document}

\maketitle
\begin{note}
this document and others can be found on my website\\
\url{https://mahmoodsheikh36.github.io/}
\end{note}
\cite{aima2010}\\
\begin{definition}
a \textbf{decision tree} represents a \href{discrete_maths2.org}{function} that takes as input a \href{linear_algebra2.org}{vector} of attribute values and returns a "decision"--a single output value. the input and output values can be \href{20221015110034-discrete_continuous_data.org}{discrete or continuous}\\
\begin{characteristic}
for a wide variety of problems, the decision tree yields a nice, concise result. but some functions cannot be represented concisely, for example, the majority function, which returns true if and only if more than half of the inputs are true, requires an exponentially large decision tree.\\
\end{characteristic}
\end{definition}
\section{inducing boolean decision trees from examples\hfill{}\textsc{code:math}}
\label{sec:org7059b94}
no matter how we measure size, it is an intractable problem to find the smallest consistent tree; there is no way to efficiently search through the \({2^2}^n\) trees. with some simple heuristics, however, we can find a good approximate solution: a small (but not smallest) consistent tree. the \(\textsc{Decision-Tree-Learning}\) algorithm adopts a \href{20230118112408-greedy_algorithm.org}{greedy} \href{20220706211939-divide_and_conquer_algorithm.org}{divide-and-conquer} strategy: always test the most important attribute first. this test divides the problem up into smaller subproblems that can then be solved recursively. by “most important attribute” we mean the one that makes the most difference to the classification of an example. that way, we hope to get to the correct classification with a small number of tests, meaning that all paths in the tree will be short and the tree as a whole will be shallow.\\
in general, after the first attribute test splits up the examples, each outcome is a new decision tree learning problem in itself, with fewer examples and one less attribute. there are four cases to consider for these recursive problems:\\
\begin{enumerate}
\item if the remaining examples are all positive (or all negative), then we are done: we can answer Yes or No\\
\item if there are some positive and some negative examples, then choose the best attribute to split them.\\
\item if there are no examples left, it means that no example has been observed for this combination of attribute values, and we return a default value calculated from the plurality classification of all the examples that were used in constructing the node’s parent. these are passed along in the variable \emph{parent\textsubscript{examples}}\\
\item if there are no attributes left, but both positive and negative examples, it means that these examples have exactly the same description, but different classifications. this can happen because there is an error or noise in the data; because the domain is nondeterministic; or because we can’t observe an attribute that would distinguish the examples. the best we can do is return the plurality classification of the remaining examples.\\
\end{enumerate}

\includesvg{/home/mahmooz/brain/out/PrkHDOF}\\

an attribute a with \(d\) distinct values divides the training set E into subsets \(E_1,\dots,E_d\). each subset \(E_k\) has \(p_k\) positive examples and \(n_k\) negative examples, so if we go along that branch, we will need an additional \(B(p_k/(p_k+n_k))\) bits of information to answer the question. a randomly chosen example from the training set has the \(k\)'th value for the attribute with probability \((p_k+n_k)/(p+n)\), so the expected \href{20230227001825-entropy.org}{entropy} remaining after testing attribute \(A\) is\\
\[
Remainder(A) = \sum_{k=1}^{d} \frac{p_k+n_k}{p+n}B\left(\frac{p_k}{p_k+n_k}\right)
\]\\
\begin{note}
this confused me at first, why multiply probability by entropy? and why would the result be entropy?\\
turns out, the probability is giving "weight" to the entropy, as the entropy of a more probable branch is more important than a less probable one\\
\end{note}
the \href{20230313145907-information_gain.org}{information gain} from the attribute test on \(A\) is the expected reduction in entropy:\\
\[
Gain(A) = B\left(\frac{p}{p+n}\right) - Remainder(A)
\]\\
which is exactly what we need to implement the \(\textsc{Importance}\) function, as the attribute that gives the most information gain is the one that should be split on\\

\lstset{language=Lisp,label= ,caption= ,captionpos=b,numbers=none}
\begin{lstlisting}
(defun decision-tree-learning (examples attributes parent-examples)
  "generate a boolean decision tree given a set of examples"
  ;; if no examples left, return the most frequent classification of the parent node's examples
  (if (not examples)
      (return-from decision-tree-learning (list 'answer (plurality-classification parent-examples))))
  ;; if all examples have the same classification, return it
  (let ((first-example (car examples))
        (same-classification t))
    (loop for example in examples
          do (if (not (equal (getf example 'goal) (getf first-example 'goal)))
                 (setf same-classification nil)))
    (if same-classification
        (return-from decision-tree-learning (list 'answer (getf first-example 'goal)))))
  ;; if no attributes left, we're done splitting, just return the most frequent classification
  (if (not attributes)
      (return-from decision-tree-learning (list 'answer (plurality-classification examples))))
  ;; find the attribute to split on, the one that gives the most information gain
  (let ((most-important-attribute (car attributes)))
    (loop for attribute in attributes
          do (let ((current-entropy (boolean-random-variable-entropy (/ (length (positive-examples examples)) (length examples)))))
               (if (> (attribute-information-gain attribute examples current-entropy) (attribute-information-gain most-important-attribute examples current-entropy))
                   (setf most-important-attribute attribute))))
    ;; create the node and its children
    (let ((tree (list 'attribute-name most-important-attribute))
          (this-attributes-possible-values (possible-attribute-values most-important-attribute examples)))
      (loop for attribute-value in this-attributes-possible-values
            do (let ((examples-with-this-attribute-value '()))
                 (loop for example in examples
                       do (if (equal (getf example most-important-attribute) attribute-value)
                              (push example examples-with-this-attribute-value)))
                 (let ((subtree (decision-tree-learning examples-with-this-attribute-value (remove most-important-attribute attributes) examples)))
                   (setf (getf subtree 'branch-value) (format nil "~a=~a" most-important-attribute attribute-value))
                   (push subtree (getf tree 'children)))))
      tree)))

(defun plurality-classification (examples)
  "get the most frequent answer, yes/no, in a set of examples"
  (let ((yes-counter 0)
        (no-counter 0))
    (loop for example in examples
          do (if (equal (getf example 'goal) "yes")
                 (incf yes-counter)
                 (incf no-counter)))
    (if (> yes-counter no-counter)
        "yes"
        "no")))

(defun possible-attribute-values (attribute examples)
  "get a list of the values that an attribute can hold"
  (let ((possible-values nil))
    (loop for example in examples
          do (pushnew (getf example attribute) possible-values :test 'equal)) ;; use equal to compare
    possible-values))

(defun boolean-random-variable-entropy (probability)
  (cond ((equal probability 1) 0) ;; always true
        ((equal probability 0) 0) ;; always false
        (t (- (+ (* probability (log probability 2)) (* (- 1 probability) (log (- 1 probability) 2)))))))

;; in pseudocode its named importance
(defun attribute-information-gain (attribute examples previous-entropy)
  "the information gain of an attribute, which reflects its importance"
  (let ((possible-values (possible-attribute-values attribute examples))
        (total-entropy 0))
    ;; each value corresponds to a theoretical child of the node (if the attribute was to be split on)
    (loop for value in possible-values
          do (let* ((examples-with-that-value (examples-with-attribute-value examples attribute value))
                    (p-examples (positive-examples examples-with-that-value))
                    (n-examples (negative-examples examples-with-that-value)))
               (incf total-entropy (* (/ (length examples-with-that-value) (length examples)) (boolean-random-variable-entropy (/ (length p-examples) (length examples-with-that-value)))))))
    (- previous-entropy total-entropy)))

(defun positive-examples (examples)
  "get the examples whose goal is yes"
  (let ((positive-examples nil))
    (loop for example in examples
          do (if (equal (getf example 'goal) "yes")
                 (push example positive-examples)))
    positive-examples))

(defun negative-examples (examples)
  "get the examples whose goal is no"
  (let ((negative-examples nil))
    (loop for example in examples
          do (if (equal (getf example 'goal) "no")
                 (push example negative-examples)))
    negative-examples))

(defun examples-with-attribute-value (examples attribute value)
  "from a set of examples, get the ones with an attribute with a specific value"
  (let ((examples-with-that-attribute-and-value nil))
    (loop for example in examples
          do (if (equal (getf example attribute) value)
                 (push example examples-with-that-attribute-and-value)))
    examples-with-that-attribute-and-value))

(defun attributes-of-example (example)
  "list of attributes of example"
  (loop for (key value) on example by #'cddr
        collect key))
\end{lstlisting}

\lstset{language=Lisp,label= ,caption= ,captionpos=b,numbers=none}
\begin{lstlisting}
(defparameter
    learning-examples
  '((alt "yes" bar "no"  fri "no"  hun "yes" pat "some" price "$$$" rain "no"  res "yes" type "french"  est  "0-10" goal "yes")
    (alt "yes" bar "no"  fri "no"  hun "yes" pat "full" price "$"   rain "no"  res "no"  type "thai"    est "30-60" goal "no" )
    (alt "no"  bar "yes" fri "no"  hun "no"  pat "some" price "$"   rain "no"  res "no"  type "burger"  est  "0-10" goal "yes")
    (alt "yes" bar "no"  fri "yes" hun "yes" pat "full" price "$"   rain "yes" res "no"  type "thai"    est "10-30" goal "yes")
    (alt "yes" bar "no"  fri "yes" hun "no"  pat "full" price "$$$" rain "no"  res "yes" type "french"  est   ">60" goal "no" )
    (alt "no"  bar "yes" fri "no"  hun "yes" pat "some" price "$$"  rain "yes" res "yes" type "italian" est  "0-10" goal "yes")
    (alt "no"  bar "yes" fri "no"  hun "no"  pat "none" price "$"   rain "yes" res "no"  type "burger"  est  "0-10" goal "no" )
    (alt "no"  bar "no"  fri "no"  hun "yes" pat "some" price "$$"  rain "yes" res "yes" type "thai"    est  "0-10" goal "yes")
    (alt "no"  bar "yes" fri "yes" hun "no"  pat "full" price "$"   rain "yes" res "no"  type "burger"  est   ">60" goal "no" )
    (alt "yes" bar "yes" fri "yes" hun "yes" pat "full" price "$$$" rain "no"  res "yes" type "italian" est "10-30" goal "no" )
    (alt "no"  bar "no"  fri "no"  hun "no"  pat "none" price "$"   rain "no"  res "no"  type "thai"    est  "0-10" goal "no" )
    (alt "yes" bar "yes" fri "yes" hun "yes" pat "full" price "$"   rain "no"  res "no"  type "burger"  est "30-60" goal "yes")))

(decision-tree-learning learning-examples (remove 'goal (attributes-of-example (car learning-examples))) nil)
\end{lstlisting}
\end{document}